\begin{proposition}[KT-SUBP-SEED-THIRD-UNIT-NONIMPLICATION]
Requiring a two-unit auxiliary discriminant to satisfy $|D|\leq p-4$
does not itself force the remaining third unit, even in the automatic
family $p\equiv3\pmod4$, $e=(p+1)/2$.
\end{proposition}
\begin{proof}
Take $p=83$, $e=42$, $b=e+1=43$, and $D=-43$.  Both $83$ and $43$
are prime, $D$ is odd and fundamental, and $|D|=43\leq79=p-4$.
The exact defining Bernoulli sums give
\[
 B_{1,\chi_{-43}}=-1,\qquad
 \beta_3(\chi_{-43})=166=2\cdot83,
 \qquad \beta_{43}(\chi_{-43})\equiv76\pmod{83}.
\]
Thus $h(-43)=1$ and the first twisted value is a unit, whereas the
third is not.  The full exact high-weight quotient is retained in the
certificate report.
The standalone certificate
\texttt{subp\_third\_unit\_counterexample.py} computes these rational
Bernoulli-polynomial values and independently obtains $76$ from the
defining conductor-$83\cdot43$ Teichm\"uller-character sum modulo $83^2$.
The field also satisfies the already checked source conditions at
$N=2,\psi=1$: its discriminant is odd, $83\nmid43h(-43)$, $e$ is
positive and even, and $82\nmid2e=84$.  The remaining quadratic and
character-intersection clauses are the same elementary verification
as in KT-JACOBI-TWO-UNIT-GD.  This is a single-field nonimplication;
it does not rule out choosing another field with all three units.
\end{proof}
